% Options for packages loaded elsewhere
\PassOptionsToPackage{unicode}{hyperref}
\PassOptionsToPackage{hyphens}{url}
\PassOptionsToPackage{dvipsnames,svgnames,x11names}{xcolor}
%
\documentclass[
  10pt,
]{article}
\usepackage{amsmath,amssymb}
\usepackage{iftex}
\ifPDFTeX
  \usepackage[T1]{fontenc}
  \usepackage[utf8]{inputenc}
  \usepackage{textcomp} % provide euro and other symbols
\else % if luatex or xetex
  \usepackage{unicode-math} % this also loads fontspec
  \defaultfontfeatures{Scale=MatchLowercase}
  \defaultfontfeatures[\rmfamily]{Ligatures=TeX,Scale=1}
\fi
\usepackage{lmodern}
\ifPDFTeX\else
  % xetex/luatex font selection
\fi
% Use upquote if available, for straight quotes in verbatim environments
\IfFileExists{upquote.sty}{\usepackage{upquote}}{}
\IfFileExists{microtype.sty}{% use microtype if available
  \usepackage[]{microtype}
  \UseMicrotypeSet[protrusion]{basicmath} % disable protrusion for tt fonts
}{}
\makeatletter
\@ifundefined{KOMAClassName}{% if non-KOMA class
  \IfFileExists{parskip.sty}{%
    \usepackage{parskip}
  }{% else
    \setlength{\parindent}{0pt}
    \setlength{\parskip}{6pt plus 2pt minus 1pt}}
}{% if KOMA class
  \KOMAoptions{parskip=half}}
\makeatother
\usepackage{xcolor}
\usepackage{color}
\usepackage{fancyvrb}
\newcommand{\VerbBar}{|}
\newcommand{\VERB}{\Verb[commandchars=\\\{\}]}
\DefineVerbatimEnvironment{Highlighting}{Verbatim}{commandchars=\\\{\}}
% Add ',fontsize=\small' for more characters per line
\newenvironment{Shaded}{}{}
\newcommand{\AlertTok}[1]{\textcolor[rgb]{1.00,0.00,0.00}{\textbf{#1}}}
\newcommand{\AnnotationTok}[1]{\textcolor[rgb]{0.38,0.63,0.69}{\textbf{\textit{#1}}}}
\newcommand{\AttributeTok}[1]{\textcolor[rgb]{0.49,0.56,0.16}{#1}}
\newcommand{\BaseNTok}[1]{\textcolor[rgb]{0.25,0.63,0.44}{#1}}
\newcommand{\BuiltInTok}[1]{\textcolor[rgb]{0.00,0.50,0.00}{#1}}
\newcommand{\CharTok}[1]{\textcolor[rgb]{0.25,0.44,0.63}{#1}}
\newcommand{\CommentTok}[1]{\textcolor[rgb]{0.38,0.63,0.69}{\textit{#1}}}
\newcommand{\CommentVarTok}[1]{\textcolor[rgb]{0.38,0.63,0.69}{\textbf{\textit{#1}}}}
\newcommand{\ConstantTok}[1]{\textcolor[rgb]{0.53,0.00,0.00}{#1}}
\newcommand{\ControlFlowTok}[1]{\textcolor[rgb]{0.00,0.44,0.13}{\textbf{#1}}}
\newcommand{\DataTypeTok}[1]{\textcolor[rgb]{0.56,0.13,0.00}{#1}}
\newcommand{\DecValTok}[1]{\textcolor[rgb]{0.25,0.63,0.44}{#1}}
\newcommand{\DocumentationTok}[1]{\textcolor[rgb]{0.73,0.13,0.13}{\textit{#1}}}
\newcommand{\ErrorTok}[1]{\textcolor[rgb]{1.00,0.00,0.00}{\textbf{#1}}}
\newcommand{\ExtensionTok}[1]{#1}
\newcommand{\FloatTok}[1]{\textcolor[rgb]{0.25,0.63,0.44}{#1}}
\newcommand{\FunctionTok}[1]{\textcolor[rgb]{0.02,0.16,0.49}{#1}}
\newcommand{\ImportTok}[1]{\textcolor[rgb]{0.00,0.50,0.00}{\textbf{#1}}}
\newcommand{\InformationTok}[1]{\textcolor[rgb]{0.38,0.63,0.69}{\textbf{\textit{#1}}}}
\newcommand{\KeywordTok}[1]{\textcolor[rgb]{0.00,0.44,0.13}{\textbf{#1}}}
\newcommand{\NormalTok}[1]{#1}
\newcommand{\OperatorTok}[1]{\textcolor[rgb]{0.40,0.40,0.40}{#1}}
\newcommand{\OtherTok}[1]{\textcolor[rgb]{0.00,0.44,0.13}{#1}}
\newcommand{\PreprocessorTok}[1]{\textcolor[rgb]{0.74,0.48,0.00}{#1}}
\newcommand{\RegionMarkerTok}[1]{#1}
\newcommand{\SpecialCharTok}[1]{\textcolor[rgb]{0.25,0.44,0.63}{#1}}
\newcommand{\SpecialStringTok}[1]{\textcolor[rgb]{0.73,0.40,0.53}{#1}}
\newcommand{\StringTok}[1]{\textcolor[rgb]{0.25,0.44,0.63}{#1}}
\newcommand{\VariableTok}[1]{\textcolor[rgb]{0.10,0.09,0.49}{#1}}
\newcommand{\VerbatimStringTok}[1]{\textcolor[rgb]{0.25,0.44,0.63}{#1}}
\newcommand{\WarningTok}[1]{\textcolor[rgb]{0.38,0.63,0.69}{\textbf{\textit{#1}}}}
\usepackage{longtable,booktabs,array}
\usepackage{calc} % for calculating minipage widths
% Correct order of tables after \paragraph or \subparagraph
\usepackage{etoolbox}
\makeatletter
\patchcmd\longtable{\par}{\if@noskipsec\mbox{}\fi\par}{}{}
\makeatother
% Allow footnotes in longtable head/foot
\IfFileExists{footnotehyper.sty}{\usepackage{footnotehyper}}{\usepackage{footnote}}
\makesavenoteenv{longtable}
\setlength{\emergencystretch}{3em} % prevent overfull lines
\providecommand{\tightlist}{%
  \setlength{\itemsep}{0pt}\setlength{\parskip}{0pt}}
\setcounter{secnumdepth}{5}

\usepackage{amsmath,amssymb}
\usepackage{graphicx}
\usepackage{longtable,booktabs,array}
\usepackage{xcolor}
\usepackage[margin=2.1cm]{geometry}
\usepackage{microtype}
\setlength{\emergencystretch}{3em}
\pdfinfoomitdate=1
\pdftrailerid{}
\pdfsuppressptexinfo=-1
\renewcommand{\arraystretch}{1.15}
\allowdisplaybreaks
\ifLuaTeX
  \usepackage{selnolig}  % disable illegal ligatures
\fi
\IfFileExists{bookmark.sty}{\usepackage{bookmark}}{\usepackage{hyperref}}
\IfFileExists{xurl.sty}{\usepackage{xurl}}{} % add URL line breaks if available
\urlstyle{same}
\hypersetup{
  pdftitle={UM-2 manual tasks},
  pdfauthor={Underlying theory: Patrick L. Nash; Written by Claude (Anthropic); Manual tasks, version 2, 2026-10-18 (version 1: 2026-10-18); Repository once-ere/Dirac\_claude, folder universe\_model\_2},
  colorlinks=true,
  linkcolor={blue},
  filecolor={Maroon},
  citecolor={Blue},
  urlcolor={blue},
  pdfcreator={LaTeX via pandoc}}

\title{UM-2 manual tasks}
\usepackage{etoolbox}
\makeatletter
\providecommand{\subtitle}[1]{% add subtitle to \maketitle
  \apptocmd{\@title}{\par {\large #1 \par}}{}{}
}
\makeatother
\subtitle{How to carry out by hand the two long tasks of the UM-2 work:
the fresh-clone verification of every computation and document, and the
independent review of the documents against the computed outputs}
\author{Underlying theory: Patrick L. Nash \and Written by Claude
(Anthropic) \and Manual tasks, version 2, 2026-10-18 (version 1:
2026-10-18) \and Repository once-ere/Dirac\_claude, folder
universe\_model\_2}
\date{}

\begin{document}
\maketitle

{
\hypersetup{linkcolor=}
\setcounter{tocdepth}{3}
\tableofcontents
}
\hypertarget{how-to-read-this-document}{%
\section{How to read this document}\label{how-to-read-this-document}}

Two tasks of the UM-2 work are long when they are run as single jobs:

\begin{itemize}
\tightlist
\item
  \textbf{Task 1, the fresh-clone verification}: all 19 computations of
  version 2 and the document build took 1027 s (about 17 minutes) on
  2026-10-18 at commit \texttt{4a6421d}, although every individual
  computation takes at most about 4.5 minutes (the primordial
  integration, 268 s).
\item
  \textbf{Task 2, the independent review}: a review agent that read the
  five model documents of version 2, compared their statements with the
  computed outputs and recomputed key numbers independently ran for 31
  minutes (1858 s, 82 tool calls).
\end{itemize}

When this document was first written the author's rule was that no
individual calculation may take longer than 12 minutes and 3 seconds
(later 7 minutes and 3 seconds); the author lifted the limit on
2026-10-18. The steps below are still kept short, so that each can be
executed by hand and its result reported back for incorporation into the
project. This document is self-contained: every command, every equation
and every reference value it needs is written here. Wherever a value can
change when the computations are updated, the command that prints it
from the repository is given instead of the value.

Version 2 (2026-10-18) covers UM-2 version 2: three generations of
quarks and leptons, three neutrino flavours with their mixing and
flavour conversion, and the dark matter of both halves of
dirac16complex. It adds review items for these (items 2.6 to 2.8) and
records the findings of the review of version 2 and their incorporation.

Conventions: the coordinates are the author's \(x_0, x_1, \dots, x_7\)
(\(x_0\) the hidden space direction, \(x_1, x_2, x_3\) space, \(x_4\)
the time, \(x_5, x_6, x_7\) the extra times); \(W^+\) carries the charge
\(+e\) and \(W^-\) the charge \(-e\).

\hypertarget{requirements-and-set-up}{%
\section{Requirements and set-up}\label{requirements-and-set-up}}

The commands are for a Linux shell: Ubuntu under WSL on Windows (start
``Ubuntu'' from the Start menu), native Linux, or the macOS terminal.
Every block can be copied and pasted as a whole.

\textbf{Step S1 (once, about 5 minutes): system packages.} This installs
git, Python, pandoc and the TeX Live parts the document build needs (on
macOS use Homebrew instead: \texttt{brew\ install\ git\ python\ pandoc}
and the MacTeX distribution).

\begin{Shaded}
\begin{Highlighting}[]
\NormalTok{sudo apt{-}get update}
\NormalTok{sudo apt{-}get install {-}y git python3 python3{-}venv python3{-}pip pandoc \textbackslash{}}
\NormalTok{    texlive{-}latex{-}base texlive{-}latex{-}recommended texlive{-}latex{-}extra \textbackslash{}}
\NormalTok{    texlive{-}fonts{-}recommended lmodern}
\end{Highlighting}
\end{Shaded}

\textbf{Step S2 (once, about 2 minutes): a Python environment with the
packages of the computations.} A virtual environment keeps these
packages separate from the system Python.

\begin{Shaded}
\begin{Highlighting}[]
\NormalTok{python3 {-}m venv \textasciitilde{}/um2env}
\NormalTok{source \textasciitilde{}/um2env/bin/activate}
\NormalTok{pip install sympy numpy scipy mpmath matplotlib camb}
\end{Highlighting}
\end{Shaded}

In every new terminal, activate it again with
\texttt{source\ \textasciitilde{}/um2env/bin/activate} before running
the commands below.

The reference runs used Python 3.13.16, sympy 1.14.0, numpy 2.5.3, scipy
1.18.1, mpmath 1.3.0, matplotlib 3.11.2, camb 2.0.5, pandoc 3.1.3 and
pdfTeX 1.40.25 (TeX Live 2023). With other versions the numbers agree to
the tolerance used below, while the PDF and LaTeX files may differ in
their bytes (not in their content).

\hypertarget{task-1-the-fresh-clone-verification}{%
\section{Task 1: the fresh-clone
verification}\label{task-1-the-fresh-clone-verification}}

\textbf{What it establishes.} Every stored number of UM-2 is reproduced
from the code alone, every named check passes, and the documents are
rebuilt from their Markdown sources.

\textbf{Step 1.1 (about 1 minute): a fresh clone, separate from any
working copy.}

\begin{Shaded}
\begin{Highlighting}[]
\NormalTok{cd \textasciitilde{}}
\NormalTok{rm {-}rf um2\_check}
\NormalTok{git clone https://github.com/once{-}ere/Dirac\_claude.git um2\_check}
\NormalTok{cd \textasciitilde{}/um2\_check/universe\_model\_2}
\NormalTok{git log {-}{-}oneline {-}1 \textgreater{} \textasciitilde{}/um2\_check\_commit.txt}
\NormalTok{cat \textasciitilde{}/um2\_check\_commit.txt}
\end{Highlighting}
\end{Shaded}

\textbf{Step 1.2 (once per terminal): a helper that runs one computation
with a time limit and records its exit code and duration.} Each
computation is stopped if it runs longer than 900 s (15 minutes, three
times the longest expected duration); the exit code is 0 if and only if
every named check of that computation passes.

\begin{Shaded}
\begin{Highlighting}[]
\NormalTok{cd \textasciitilde{}/um2\_check/universe\_model\_2}
\NormalTok{run() \{}
\NormalTok{  local t0=$(date +\%s)}
\NormalTok{  timeout {-}k 5 900 python3 "$@" \textgreater{} \textasciitilde{}/um2\_log\_$(echo "$*" | tr \textquotesingle{} /\textquotesingle{} \textquotesingle{}\_\_\textquotesingle{}).txt 2\textgreater{}\&1}
\NormalTok{  local rc=$?}
\NormalTok{  echo "$* exit=$rc $(( $(date +\%s) {-} t0 )) s" | tee {-}a \textasciitilde{}/um2\_times.txt}
\NormalTok{\}}
\NormalTok{: \textgreater{} \textasciitilde{}/um2\_times.txt}
\end{Highlighting}
\end{Shaded}

\textbf{Step 1.3: the computations, one block at a time, in this order}
(each block depends on the outputs of the blocks before it; the expected
durations are those of the run of 2026-10-18 at commit
\texttt{4a6421d}).

\begin{longtable}[]{@{}
  >{\raggedright\arraybackslash}p{(\columnwidth - 6\tabcolsep) * \real{0.0541}}
  >{\raggedright\arraybackslash}p{(\columnwidth - 6\tabcolsep) * \real{0.6757}}
  >{\raggedright\arraybackslash}p{(\columnwidth - 6\tabcolsep) * \real{0.1622}}
  >{\raggedright\arraybackslash}p{(\columnwidth - 6\tabcolsep) * \real{0.1081}}@{}}
\toprule\noalign{}
\begin{minipage}[b]{\linewidth}\raggedright
block
\end{minipage} & \begin{minipage}[b]{\linewidth}\raggedright
commands
\end{minipage} & \begin{minipage}[b]{\linewidth}\raggedright
expected time
\end{minipage} & \begin{minipage}[b]{\linewidth}\raggedright
checks
\end{minipage} \\
\midrule\noalign{}
\endhead
\bottomrule\noalign{}
\endlastfoot
1 & \texttt{run\ code/um2\_algebra.py} & 54 s & 31 \\
2 & \texttt{run\ code/um2\_geometry.py} & 77 s & 25 \\
3 & \texttt{run\ code/um2\_background.py} & 133 s & 22 \\
4 & \texttt{run\ code/um2\_fields.py} & 54 s & 16 \\
5 & \texttt{run\ code/um2\_spectra.py} & 28 s & 14 \\
6 & \texttt{run\ code/um2\_flavour.py} and
\texttt{run\ code/um2\_neutrino.py} & 5 s, 31 s & 7, 15 \\
7 & \texttt{run\ code/um2\_ks\_jobs.py\ validate} and
\texttt{run\ code/um2\_ks\_jobs.py\ spectrum} & 37 s, 1 s & 2, 3 \\
8 & \texttt{run\ code/um2\_ks\_jobs.py\ ground8} & 137 s & 3 \\
9 & \texttt{run\ code/um2\_ks\_jobs.py\ ground592} & 137 s & 4 \\
10 & \texttt{run\ code/um2\_ks\_jobs.py\ darkmatter},
\texttt{run\ code/um2\_ks\_jobs.py\ portals},
\texttt{run\ code/um2\_ks\_jobs.py\ feedback} & 16 s, 4 s, 0.5 s & 3, 5,
2 \\
11 & \texttt{run\ code/um2\_electroweak.py} & 1 s & 19 \\
12 & \texttt{run\ code/um2\_primordial.py} & 268 s & 5 \\
13 & \texttt{run\ code/um2\_cosmology.py} and
\texttt{run\ code/um2\_analysis.py} & 23 s, 4 s & 6, 23 \\
14 & \texttt{run\ tools/build\_docs.py} & 17 s & \\
\end{longtable}

The check counts add up to 205; the counts in force are printed by step
1.4. For example, block 1:

\begin{Shaded}
\begin{Highlighting}[]
\NormalTok{run code/um2\_algebra.py}
\end{Highlighting}
\end{Shaded}

\textbf{Step 1.4 (seconds): the summary of the checks.} Every report
must say \texttt{0\ failed}; the last line gives the total.

\begin{Shaded}
\begin{Highlighting}[]
\NormalTok{grep {-}h "checks, " \textasciitilde{}/um2\_log\_*.txt}
\NormalTok{grep {-}h "checks, " \textasciitilde{}/um2\_log\_*.txt | awk \textquotesingle{}\{s += $2; f += $4\} END \{print "TOTAL", s, "checks,", f, "failed"\}\textquotesingle{}}
\NormalTok{cat \textasciitilde{}/um2\_times.txt}
\end{Highlighting}
\end{Shaded}

\textbf{Step 1.5 (seconds): the comparison with the committed files.}
The script lists every file the rerun changed and compares the numbers
of every changed JSON file with the committed values to a relative
tolerance of \(10^{-9}\).

\begin{Shaded}
\begin{Highlighting}[]
\NormalTok{cd \textasciitilde{}/um2\_check}
\NormalTok{python3 universe\_model\_2/tools/compare\_outputs.py}
\end{Highlighting}
\end{Shaded}

Expected result: \texttt{result:\ REPRODUCED}. In the run of 2026-10-18
at commit \texttt{4a6421d} only the run times and one Kohn-Sham integral
(a relative difference of \(1.1\times10^{-12}\), caused by threaded
linear algebra that does not sum in a fixed order) differed; the LaTeX
sources, the PDF documents and the figures were byte-identical. With
other software versions, files reported as \texttt{CHANGED\ (not\ JSON)}
(PDF, LaTeX, PNG) are expected and are not errors; open them to confirm
that their content is the same.

\textbf{Step 1.6: what to report back.} The contents of
\texttt{\textasciitilde{}/um2\_check\_commit.txt},
\texttt{\textasciitilde{}/um2\_times.txt}, the output of steps 1.4 and
1.5, and, for any computation with a non-zero exit code, the last 40
lines of its log
(\texttt{tail\ -40\ \textasciitilde{}/um2\_log\_\textless{}name\textgreater{}.txt}).

\hypertarget{task-2-the-independent-review-of-the-documents}{%
\section{Task 2: the independent review of the
documents}\label{task-2-the-independent-review-of-the-documents}}

\textbf{What it establishes.} That the statements, equations and numbers
of the five model documents (MODEL, PROVENANCE, FIELD\_EQUATIONS,
SOLUTION, ANALYSIS) agree with the computed outputs and with each other,
that the mathematics written in them is correct, and that each document
is self-contained (every equation it uses is written in it, and it never
sends the reader to another document for content). Each item takes 5 to
15 minutes by hand. Work in the fresh clone of task 1 (step 1.1) or in
any up-to-date clone:

\begin{Shaded}
\begin{Highlighting}[]
\NormalTok{cd \textasciitilde{}/um2\_check/universe\_model\_2}
\end{Highlighting}
\end{Shaded}

\textbf{Item 2.1: the electroweak parameter, the \(W^\pm\), \(Z\) fit
and the \(Z\) widths.} The tip term of the sigma field gives the \(x_0\)
eigenvalue problem of the \(W^\pm\) and \(Z\) profiles the tip
coefficients \(\beta_W\) and \(\beta_Z = \beta_W/\cos^2\theta_W\), with
\(\theta_W\) the angle of the gauge couplings; the eigenvalues
\(\Lambda_W\), \(\Lambda_Z\) are the masses in units of the proper half
cell, \(\Lambda = (m\,fL)^2\), and
\[ \rho = \frac{\Lambda_W}{\Lambda_Z\cos^2\theta_W},\qquad \cos^2\theta_W = \frac{m_W^2/m_Z^2}{\rho_0},\qquad m_{KK} = m_W\sqrt{\Lambda_1/\Lambda_W},\qquad m_Z^{model} = m_W\sqrt{\Lambda_Z/\Lambda_W}, \]
with \(\Lambda_1\) the first even Kaluza-Klein eigenvalue and
\(\rho_0 - 1 = 0.00038\) the measured value. The tree-level partial
widths are
\(\Gamma(Z\to f\bar f) = N_c\,G_Fm_Z^3(g_V^2 + g_A^2)/(6\sqrt2\pi)\),
\(g_V = T_3 - 2Q\sin^2\theta_W\), \(g_A = T_3\), summed over the three
generations (the top quark excluded). Print what the code computes:

\begin{Shaded}
\begin{Highlighting}[]
\NormalTok{python3 {-}c "import json; print(json.load(open(\textquotesingle{}outputs/masses.json\textquotesingle{}))[\textquotesingle{}WZ\textquotesingle{}])"}
\NormalTok{python3 {-}c "import json; e = json.load(open(\textquotesingle{}outputs/electroweak.json\textquotesingle{})); print(\{k: e[k] for k in e if \textquotesingle{}width\textquotesingle{} in k.lower() or \textquotesingle{}N\_nu\textquotesingle{} in k\})"}
\NormalTok{grep {-}n "c2 = \textbackslash{}|rho\_minus\_1\textbackslash{}|def LamW" code/um2\_spectra.py}
\NormalTok{grep {-}n "rho = \textbackslash{}|cos\^{}2\textbackslash{}\textbackslash{}\textbackslash{}\textbackslash{}theta\_W\textbackslash{}|m\_Z" MODEL.md SOLUTION.md ANALYSIS.md FIELD\_EQUATIONS.md PROVENANCE.md}
\end{Highlighting}
\end{Shaded}

Check: (a) the angle the code uses (the variable \texttt{c2} in
\texttt{code/um2\_spectra.py}) is the angle the documents name; (b) the
documents' \(\beta_W\), \(\beta_Z\), \(\Lambda_W\), \(\Lambda_Z\),
\(\rho - 1\), \(m_{KK}\) and \(m_Z^{model}\) equal the printed values
(rounded), and \(m_Z^{model}\) equals the measured \(91.1879\) GeV; (c)
the \(Z\) widths of the documents (invisible about 498 MeV, total about
2485 MeV with the factor \(1 + \alpha_s/\pi\) on the quarks) equal the
printed values.

\textbf{Item 2.2: the uncertainty relation on the cell.} With the scaled
hidden coordinate \(\xi = y/L \in [-1, 0]\) (fold at \(\xi = 0\), tip at
\(\xi = -1\)), a normalised probability density \(\varrho\)
(\(\int_{-1}^{0}\varrho\,d\xi = 1\)), the spread
\(\Delta\xi^2 = \int_{-1}^{0}\xi^2\varrho\,d\xi\) and the hidden kinetic
term \(\langle p^2\rangle\) (bosons:
\(\varrho = \tilde u^2/\!\int\tilde u^2\) with
\(\tilde u = e^{3\eta\xi/2}u\) and
\(\langle p^2\rangle = \int\tilde u'^2/\!\int\tilde u^2\); fermions:
\(\varrho = |\chi|^2/\!\int|\chi|^2\) and
\(\langle p^2\rangle = \int|\chi'|^2/\!\int|\chi|^2\); primes
\(d/d\xi\)), the documents state
\[ \sqrt{\langle p^2\rangle}\,\Delta\xi \ge \frac{1 - \varrho(\text{tip})}{2}. \]
Verify the derivation by hand: integration by parts gives
\(1 = \int_{-1}^{0}\varrho\,d\xi = [\xi\varrho]_{-1}^{0} - \int_{-1}^{0}\xi\varrho'\,d\xi = \varrho(-1) - \int_{-1}^{0}\xi\varrho'\,d\xi\);
with \(\varrho' = 2\tilde u\tilde u'/\!\int\tilde u^2\) (or
\(2\,\mathrm{Re}\,\chi^\dagger\chi'/\!\int|\chi|^2\)) the Cauchy-Schwarz
inequality gives
\(\lvert\int\xi\varrho'\,d\xi\rvert \le 2\Delta\xi\sqrt{\langle p^2\rangle}\).
Also verify the boson Rayleigh quotient
\((m\,fL)^2 = \big[\int e^{3\eta\xi}u'^2d\xi + \beta e^{-3\eta}u(-1)^2\big]/\!\int e^{3\eta\xi}u^2d\xi\)
and its decomposition
\((m\,fL)^2 = \langle p^2\rangle + \tfrac{9\eta^2}{4} - \tfrac{3\eta}{2}\varrho(0) + (\tfrac{3\eta}{2} + \beta)\varrho(-1)\),
which follows from
\(e^{3\eta\xi}u'^2 = \tilde u'^2 + \tfrac94\eta^2\tilde u^2 - \tfrac32\eta(\tilde u^2)'\).
Then compare the numbers:

\begin{Shaded}
\begin{Highlighting}[]
\NormalTok{python3 {-} \textless{}\textless{}\textquotesingle{}EOF\textquotesingle{}}
\NormalTok{import json}
\NormalTok{d = json.load(open(\textquotesingle{}outputs/analysis.json\textquotesingle{}))[\textquotesingle{}decomposition\textquotesingle{}]}
\NormalTok{for k, v in d.items():}
\NormalTok{    print(\textquotesingle{}\%{-}45s rho\_tip \%.4g  \textless{}p2\textgreater{}Dxi2 \%.4g  lhs \%.4f  bound \%.4f\textquotesingle{} \% (}
\NormalTok{        k, v[\textquotesingle{}rho\_tip\textquotesingle{}], v[\textquotesingle{}p2\_Dt2\textquotesingle{}], v[\textquotesingle{}uncertainty\_lhs\textquotesingle{}], v[\textquotesingle{}uncertainty\_bound\textquotesingle{}]))}
\NormalTok{EOF}
\NormalTok{python3 {-}c "import json; [print(c[\textquotesingle{}name\textquotesingle{}], c[\textquotesingle{}passed\textquotesingle{}]) for c in json.load(open(\textquotesingle{}reports/analysis{-}checks.json\textquotesingle{}))[\textquotesingle{}checks\textquotesingle{}] if \textquotesingle{}uncertainty\textquotesingle{} in c[\textquotesingle{}name\textquotesingle{}]]"}
\end{Highlighting}
\end{Shaded}

Check: every state satisfies \texttt{lhs\ \textgreater{}=\ bound}; the
tables of the ANALYSIS section on the uncertainty relation (columns
\(\varrho(\text{tip})\) and \(\langle p^2\rangle\Delta\xi^2\)) show
these values; the documents do not claim that
\(\langle p^2\rangle \ge 1/(4\Delta\xi^2)\) holds for every state (the
broad fold zero mode of dirac16complex violates it).

\textbf{Item 2.3: notation, coordinates and cross-references.} Search
for leftovers of earlier notation and for relabelled coordinates (each
command should print nothing, or only lines where the symbol has the
meaning given in the document's list of symbols with more than one use):

\begin{Shaded}
\begin{Highlighting}[]
\NormalTok{grep {-}n \textquotesingle{}t = y/L\textbackslash{}|\textbackslash{}\textbackslash{}eta t\textbackslash{}|Delta t\textbackslash{}b\textquotesingle{} *.md}
\NormalTok{grep {-}n \textquotesingle{}d\textbackslash{}\textbackslash{}psi/dt\textbackslash{}|/dt\textbackslash{}b\textbackslash{}|\textbackslash{}\textbackslash{}Omega t\textbackslash{}|Kz\textbackslash{}b\textquotesingle{} *.md}
\NormalTok{grep {-}n \textquotesingle{}r\_3(n\^{}2)\textquotesingle{} *.md}
\NormalTok{grep {-}n \textquotesingle{}\^{}\#\textquotesingle{} MODEL.md FIELD\_EQUATIONS.md SOLUTION.md ANALYSIS.md}
\end{Highlighting}
\end{Shaded}

Check: the time is always written \(x_4\) (never \(t\)) and the symmetry
axis of the fast-flavour analysis \(x_3\) (never \(z\), which means
\(6Hx_0\)); each document has a list of symbols with more than one use
(near its beginning) whose claims are true for that document; the
section numbers cited in the documents match the numbering of the last
command (pandoc numbers the sections in order, ``How to read this
document'' being section 1); no document sends the reader to another
document for an equation or a value.

\textbf{Item 2.4: the cell volume.} On the static block the metric is
\(ds^2 = e^{2a_4}d\vec x^{\,2} - dx_4^2 + f^2\big[dy^2 - e^{2Hy}d\vec x_t^{\,2}\big]\),
so the proper volume element of a slice \(x_4 = \text{const}\) is
\(e^{3a_4}f^4e^{3Hy}\,dy\,d^3x\,d^3x_t\). With
\(\mathrm{Vol}_7 = (2\pi/\Delta k)^3v_tf^4\) (at \(a_4 = 0\); \(v_t\)
the coordinate volume of the extra times) the proper seven-volume of one
cell (two halves) is
\(2\,\mathrm{Vol}_7\int_{-L}^{0}e^{3Hy}dy = 2\,\mathrm{Vol}_7(1 - e^{-3\eta})/(3H)\),
\(\eta = HL\), and the transverse volume of the four-dimensional
reduction is \(V_\perp = 2f^4v_t(1 - e^{-3\eta})/(3H)\). Check these
statements in the Kohn-Sham sections of MODEL, FIELD\_EQUATIONS and
SOLUTION, and the number \(\mathrm{Vol}_7 = 34.27\)
(\(\Delta k = 0.25H\), \(v_t = 1\), \(H = 1\)):

\begin{Shaded}
\begin{Highlighting}[]
\NormalTok{python3 {-}c "import json, math; f = json.load(open(\textquotesingle{}outputs/background.json\textquotesingle{}))[\textquotesingle{}casimir\_type\textquotesingle{}][\textquotesingle{}f\_star\textquotesingle{}]; print(f, (2*math.pi/0.25)**3 * f**4)"}
\end{Highlighting}
\end{Shaded}

\textbf{Item 2.5: the CMB statements.} The relative deviation of a
spectrum is \(\delta_\ell = C_\ell/C_\ell^{\Lambda CDM} - 1\), the
cosmic variance of one multipole \(\sqrt{2/(2\ell + 1)}\). The stored
models are, in this order, \(\Lambda\)CDM (CAMB defaults), UM-2 (three
massive neutrinos, no dark radiation) and two variants with a
thermalised fold band of dirac16complex (decoupled above the electroweak
scale and above the QCD crossover). Print the stored values:

\begin{Shaded}
\begin{Highlighting}[]
\NormalTok{python3 {-} \textless{}\textless{}\textquotesingle{}EOF\textquotesingle{}}
\NormalTok{import json}
\NormalTok{c = json.load(open(\textquotesingle{}outputs/cosmology.json\textquotesingle{}))}
\NormalTok{for k in (\textquotesingle{}max\_rel\_dev\_TT\textquotesingle{}, \textquotesingle{}max\_rel\_dev\_EE\textquotesingle{}, \textquotesingle{}max\_dev\_in\_cosmic\_variance\_units\textquotesingle{}, \textquotesingle{}UM2\_max\_rel\_dev\_TT\textquotesingle{},}
\NormalTok{          \textquotesingle{}dark\_radiation\_QCD\_max\_rel\_dev\_TT\textquotesingle{}, \textquotesingle{}dark\_radiation\_QCD\_max\_dev\_in\_cosmic\_variance\_units\textquotesingle{},}
\NormalTok{          \textquotesingle{}dark\_radiation\_QCD\_TT\_chi2\_cosmic\_variance\_full\_sky\textquotesingle{}):}
\NormalTok{    print(k, c[k])}
\NormalTok{lcdm, um2, var\_ew, var\_qcd = c[\textquotesingle{}models\textquotesingle{}]}
\NormalTok{for m in c[\textquotesingle{}models\textquotesingle{}]:}
\NormalTok{    print(m[\textquotesingle{}name\textquotesingle{}], \textquotesingle{} BAO chi2 =\textquotesingle{}, m[\textquotesingle{}bao\_chi2\textquotesingle{}], \textquotesingle{} H0 =\textquotesingle{}, m[\textquotesingle{}H0\textquotesingle{}])}
\NormalTok{for (l, a), (\_, b) in zip(lcdm[\textquotesingle{}cl\_TT\_l\textquotesingle{}], var\_qcd[\textquotesingle{}cl\_TT\_l\textquotesingle{}]):}
\NormalTok{    print(\textquotesingle{}l = \%4d  QCD variant/LCDM {-} 1 = \%+.4f   cosmic variance \%.4f\textquotesingle{} \% (l, b / a {-} 1, (2 / (2 * l + 1)) ** 0.5))}
\NormalTok{EOF}
\end{Highlighting}
\end{Shaded}

Check: the maxima are taken over \(30 \le \ell \le 2500\) (see
\texttt{code/um2\_cosmology.py}, the line
\texttt{sel\ =\ (l\ \textgreater{}=\ 30)\ \&\ (l\ \textless{}=\ 2500)}),
and the documents state that range; the cosmic-variance fraction quoted
for UM-2 refers to TT (the code computes it for TT); the BAO \(\chi^2\)
of the documents (UM-2 31.72 for 13 points, \(\Lambda\)CDM 32.55) equal
the printed values.

\textbf{Item 2.6: three generations, masses and mixing.} With one
doublet localisation per sector (\(c_L\) for the leptons, \(c_Q\) for
the quarks) and a unitary tip matrix \(\kappa\), the mass of each
fermion is the root of
\[ D(e) = e^2s_Ls_R - (C_L - c_Ls_L)(C_R + c_Rs_R) = 0,\qquad s = \frac{\sinh w}{w},\quad C = \cosh w,\quad w = \sqrt{c^2 - e^2},\qquad e = m\,fL, \]
whose first-order solution is \(e \simeq \zeta(c_L)\zeta(-c_R)\),
\(\zeta(c)^2 = c\,e^{-c}/\sinh c\); the mixing matrices are
\(U = \kappa_e^\dagger\kappa_\nu\) and \(V = \kappa_u^\dagger\kappa_d\),
equal to \(\kappa_\nu\) and \(\kappa_d\) only with
\(\kappa_e = \kappa_u = 1\) (ASSUMED). Verify by hand that for
\(c \gg 1\) \(\zeta(c)^2 \simeq 2c\,e^{-2c}\), and print the stored
values:

\begin{Shaded}
\begin{Highlighting}[]
\NormalTok{python3 {-} \textless{}\textless{}\textquotesingle{}EOF\textquotesingle{}}
\NormalTok{import json}
\NormalTok{m = json.load(open(\textquotesingle{}outputs/masses.json\textquotesingle{}))[\textquotesingle{}fermions\textquotesingle{}]}
\NormalTok{for k, v in m.items():}
\NormalTok{    print(\textquotesingle{}\%{-}4s m = \%.6g GeV  c\_L = \%.4f  c\_R = \%.4f  e\_exact = \%.6g  e\_first\_order = \%.6g\textquotesingle{} \% (}
\NormalTok{        k, v[\textquotesingle{}m\_GeV\textquotesingle{}], v[\textquotesingle{}c\_L\textquotesingle{}], v[\textquotesingle{}c\_R\textquotesingle{}], v[\textquotesingle{}e\_exact\textquotesingle{}], v[\textquotesingle{}e\_first\_order\textquotesingle{}]))}
\NormalTok{f = json.load(open(\textquotesingle{}outputs/flavour.json\textquotesingle{}))}
\NormalTok{print(\textquotesingle{}PMNS angles\textquotesingle{}, f[\textquotesingle{}PMNS\textquotesingle{}][\textquotesingle{}angles\_deg\textquotesingle{}], \textquotesingle{} J =\textquotesingle{}, f[\textquotesingle{}PMNS\textquotesingle{}][\textquotesingle{}J\textquotesingle{}], \textquotesingle{} J\_max =\textquotesingle{}, f[\textquotesingle{}PMNS\textquotesingle{}][\textquotesingle{}J\_max\textquotesingle{}])}
\NormalTok{print(\textquotesingle{}CKM J =\textquotesingle{}, f[\textquotesingle{}CKM\textquotesingle{}][\textquotesingle{}J\textquotesingle{}])}
\NormalTok{for k, v in f[\textquotesingle{}three\_generation\_tip\_condition\textquotesingle{}].items():}
\NormalTok{    print(k, \textquotesingle{}null vectors vs kappa columns:\textquotesingle{}, v[\textquotesingle{}null\_vector\_deviation\_from\_kappa\_columns\textquotesingle{}])}
\NormalTok{EOF}
\end{Highlighting}
\end{Shaded}

Check: the twelve masses, localisations and the mixing values of the
documents equal the printed ones; the documents say which \(J_{max}\) is
computed from the best-fit angles (about \(0.0335\)) and which is
NuFIT's marginalised value (\(0.0333 \pm 0.0007\)); the three-generation
determinant is said to reproduce the one-generation roots (to
\(3.0\times10^{-42}\) in the 50-digit computation).

\textbf{Item 2.7: the flavour conversion of the three neutrinos.} The
flavour amplitudes of a neutrino of energy \(E\) evolve with
\(i\,d\psi/dx_4 = [U\,\mathrm{diag}(m_1^2, m_2^2, m_3^2)U^\dagger/(2E) + \mathrm{diag}(V_{CC}, 0, 0)]\psi\),
\(V_{CC} = \sqrt2G_Fn_e\). Verify by hand: (a) the vacuum probability
\(P(\nu_\alpha\to\nu_\beta) = \lvert\sum_iU_{\beta i}U^*_{\alpha i}e^{-im_i^2L/(2E)}\rvert^2\)
and the T-odd difference
\(P(\nu_\mu\to\nu_e) - P(\nu_e\to\nu_\mu) = -16J\sin\Delta_{21}\sin\Delta_{31}\sin\Delta_{32}\),
\(\Delta_{ij} = \Delta m^2_{ij}L/(4E)\); (b) the 1-3 resonance condition
\(2EV_{CC} = \Delta m^2_{ee}\cos2\theta_{13}\),
\(\Delta m^2_{ee} = c_{12}^2\Delta m^2_{31} + s_{12}^2\Delta m^2_{32}\)
(about 11.1 GeV at 2.8 g/cm\(^3\), \(Y_e = 0.5\)); (c) the adiabatic
solar survival probability
\(P_{ee} = c_{13}^4(\frac12 + \frac12\cos2\theta_{12}\cos2\theta^m_{12}) + s_{13}^4\)
with
\(\cos2\theta^m_{12} = (\cos2\theta_{12} - A)/\sqrt{(\cos2\theta_{12} - A)^2 + \sin^22\theta_{12}}\),
\(A = 2EV_{CC}c_{13}^2/\Delta m^2_{21}\); (d) the CNB density
\(n_i = 3\zeta(3)T_\nu^3/(2\pi^2)\) per mass state with
\(T_\nu = (4/11)^{1/3}T_{CMB}\) (about 112 cm\(^{-3}\)). Print the
stored values:

\begin{Shaded}
\begin{Highlighting}[]
\NormalTok{python3 {-} \textless{}\textless{}\textquotesingle{}EOF\textquotesingle{}}
\NormalTok{import json}
\NormalTok{n = json.load(open(\textquotesingle{}outputs/neutrino.json\textquotesingle{}))}
\NormalTok{print(\textquotesingle{}masses (eV)\textquotesingle{}, n[\textquotesingle{}masses\_eV\textquotesingle{}], \textquotesingle{} sum\textquotesingle{}, n[\textquotesingle{}sum\_masses\_eV\textquotesingle{}])}
\NormalTok{print(\textquotesingle{}MSW\textquotesingle{}, n[\textquotesingle{}MSW\_resonance\textquotesingle{}])}
\NormalTok{for k, v in n[\textquotesingle{}long\_baseline\textquotesingle{}].items():}
\NormalTok{    print(k, \textquotesingle{}P(mu{-}\textgreater{}e) =\textquotesingle{}, v[\textquotesingle{}P\_mue\_nu\textquotesingle{}], \textquotesingle{} Pbar =\textquotesingle{}, v[\textquotesingle{}P\_mue\_antinu\textquotesingle{}], \textquotesingle{} vacuum =\textquotesingle{}, v[\textquotesingle{}P\_mue\_nu\_vacuum\textquotesingle{}])}
\NormalTok{for s in n[\textquotesingle{}solar\textquotesingle{}]:}
\NormalTok{    print(\textquotesingle{}solar\textquotesingle{}, s)}
\NormalTok{print(\textquotesingle{}onset T (MeV)\textquotesingle{}, n[\textquotesingle{}oscillation\_onset\_T\_MeV\textquotesingle{}])}
\NormalTok{print(\textquotesingle{}N\_eff computed\textquotesingle{}, n[\textquotesingle{}N\_eff\_UM2\_computed\textquotesingle{}], \textquotesingle{} adopted\textquotesingle{}, n[\textquotesingle{}N\_eff\_adopted\textquotesingle{}][\textquotesingle{}value\textquotesingle{}])}
\NormalTok{print(\textquotesingle{}CNB\textquotesingle{}, \{k: n[\textquotesingle{}CNB\textquotesingle{}][k] for k in (\textquotesingle{}T\_nu\_K\textquotesingle{}, \textquotesingle{}n\_per\_mass\_state\_cm3\textquotesingle{}, \textquotesingle{}Omega\_nu\_h2\textquotesingle{})\})}
\NormalTok{print(\textquotesingle{}supernova\textquotesingle{}, n[\textquotesingle{}supernova\_MSW\textquotesingle{}][\textquotesingle{}P\_ee\_normal\_ordering\textquotesingle{}], n[\textquotesingle{}supernova\_MSW\textquotesingle{}][\textquotesingle{}Pbar\_ee\_normal\_ordering\textquotesingle{}])}
\NormalTok{for k, v in n[\textquotesingle{}fast\_flavour\_instability\textquotesingle{}].items():}
\NormalTok{    print(k, \textquotesingle{} max growth\textquotesingle{}, v[\textquotesingle{}max\_growth\_any\_k\textquotesingle{}], \textquotesingle{} at k\textquotesingle{}, v[\textquotesingle{}k\_at\_max\textquotesingle{}])}
\NormalTok{EOF}
\end{Highlighting}
\end{Shaded}

Check: every number of the flavour-conversion sections of the documents
equals the printed value (rounded); the adiabatic and numerical solar
values are distinguished; the fast-flavour dispersion relation is
written with \(\omega = \Omega - \Lambda^0\), \(k = K - \Lambda^3\)
(upper indices) and \(\eta^{\mu\nu} = \mathrm{diag}(1, -1)\) on
\((x_4, x_3)\).

\textbf{Item 2.8: the dark-matter relic and unitarity.} The operator
\(G\,\bar\chi\chi F_{\mu\nu}F^{\mu\nu}\) annihilates a Dirac pair in the
p-wave, \(\sigma v = G^2s(s - 4M^2)/(4\pi) \to G^2M^4v^2/\pi\), and only
from the state \(^3P_0\) (\(J = 0\)); partial-wave unitarity averaged
over the four spin states then bounds \(\sigma \le \pi/(4p^2)\),
\(p = Mv/2\), i.e.~\(\sigma v \le \pi/(M^2v)\), whose thermal average
over \(f(v) \propto v^2e^{-xv^2/4}\) is
\[ \Big\langle\frac{\pi}{M^2v}\Big\rangle = \frac{\sqrt{\pi x}}{M^2} . \]
Verify the average by hand (the integrals
\(\int_0^\infty v\,e^{-xv^2/4}dv = 2/x\) and
\(\int_0^\infty v^2e^{-xv^2/4}dv = 2\sqrt\pi\,x^{-3/2}\), with
\(1\ \text{GeV}^{-2} = 1.1673\times10^{-17}\) cm\(^3\)/s for
\(\sigma v\)), and print the stored values:

\begin{Shaded}
\begin{Highlighting}[]
\NormalTok{python3 {-}c "import json; print(json.load(open(\textquotesingle{}outputs/analysis.json\textquotesingle{}))[\textquotesingle{}relic\textquotesingle{}])"}
\NormalTok{python3 {-}c "import json; [print(c[\textquotesingle{}detail\textquotesingle{}]) for c in json.load(open(\textquotesingle{}reports/analysis{-}checks.json\textquotesingle{}))[\textquotesingle{}checks\textquotesingle{}] if c[\textquotesingle{}name\textquotesingle{}].startswith(\textquotesingle{}B3\textquotesingle{})]"}
\NormalTok{grep {-}n "unitarity" MODEL.md SOLUTION.md ANALYSIS.md PROVENANCE.md}
\end{Highlighting}
\end{Shaded}

Check: the documents state that the freeze-out requirement (about
\(2.0\times10^{-25}\) cm\(^3\)/s per species) exceeds the bound (about
\(5.8\times10^{-26}\) cm\(^3\)/s at \(x_f = 30.4\), \(M = 44.35\) TeV)
about 3.5 times, so that the relic abundance is OPEN in version 2, and
that no document calls the candidate ``viable'' without this condition.

\textbf{Item 2.9: the built files.} Rebuild and inspect:

\begin{Shaded}
\begin{Highlighting}[]
\NormalTok{python3 tools/build\_docs.py}
\NormalTok{grep {-}c \textquotesingle{}\textbackslash{}\textbackslash{}\textbackslash{}$\textquotesingle{} MODEL.tex PROVENANCE.tex FIELD\_EQUATIONS.tex SOLUTION.tex ANALYSIS.tex}
\NormalTok{grep {-}n \textquotesingle{}textbackslash\textquotesingle{} MODEL.tex PROVENANCE.tex FIELD\_EQUATIONS.tex SOLUTION.tex ANALYSIS.tex}
\end{Highlighting}
\end{Shaded}

Check: every document reports \texttt{overfull\ boxes:\ 0}; there is no
\texttt{\textbackslash{}\$} (an unparsed dollar sign) in the LaTeX files
of the five model documents (this document contains it once, in the text
of this check); \texttt{textbackslash} appears only inside
\texttt{\textbackslash{}texorpdfstring\{...\}\{...\}} (the plain-text
PDF bookmarks of section titles).

\textbf{What to report back.} For each item: OK, or the file, the line,
what is wrong and the proposed correction.

\hypertarget{the-findings-of-the-review-of-version-1-and-their-incorporation}{%
\section{The findings of the review of version 1 and their
incorporation}\label{the-findings-of-the-review-of-version-1-and-their-incorporation}}

The review agent reported four errors and ten minor items for the
documents of version 1 at commit \texttt{9285bb6}:

\begin{enumerate}
\def\labelenumi{\arabic{enumi}.}
\tightlist
\item
  \textbf{The familiar uncertainty bound.} The documents said that
  \(\langle p^2\rangle \ge 1/(4\Delta\xi^2)\) ``holds only'' for
  profiles negligible at the tip; states with a large tip density can
  satisfy it too (first even Kaluza-Klein mode:
  \(\varrho(\text{tip}) = 0.655\),
  \(\langle p^2\rangle\Delta\xi^2 = 4.59\)). Correct wording: the bound
  is guaranteed only for such profiles.
\item
  \textbf{The symbol list of PROVENANCE} said that \(\rho_0\) is the
  electroweak parameter (PROVENANCE uses \(\rho\)) and named the lattice
  vector \(\vec n\), which PROVENANCE does not use, while \(n\) is used
  there as a level index and as the basis size of the Richardson
  formula.
\item
  \textbf{Primes.} Lists stating that primes denote \(d/d\xi\)
  overlooked \(G'(k_*)\) and \(U'(S)\) (ANALYSIS) and the reduced
  actions, whose \(u'\) is \(du/dy\) (FIELD\_EQUATIONS and MODEL).
\item
  \textbf{The electroweak fit.} With the angle of the couplings set
  equal to the on-shell angle of the measured masses,
  \(\cos^2\theta_W = m_W^2/m_Z^2\), a fit of \(\rho - 1\) to the
  measured \(\rho_0 - 1 = 0.00038\) makes the model's \(m_Z = 91.1706\)
  GeV, 17 MeV below the measured \(91.1879 \pm 0.0020\) GeV; the
  documents did not say so. (\(\rho_0\) measures the custodial violation
  beyond the Standard-Model loops; in a tree-level model it is the
  tree-level \(\rho\) with the angle of the couplings.)
\end{enumerate}

Minor items: the probability density was defined twice in ANALYSIS
(unnormalised in one section); the largest TT deviation of a
thermalised-fold-band variant lies at the edge \(\ell = 2500\) of the
range; the cosmic-variance fraction refers to TT only; gaps in the lists
of symbols; \(\mathrm{Vol}_7\) not defined in PROVENANCE; the
verification statements credited an older commit; ``true lower bound''
wording for \(m \ge m_H\), which is observed for the computed states,
not derived; the sign convention of \(w_a\); old notation in the report
texts of two checks; ``every state'' meaning every tabulated state.

All four errors and the minor items were incorporated on 2026-10-18. For
error 4 the computation was changed, not only the text: the angle of the
couplings is now fixed by the measured masses and \(\rho_0\),
\(\cos^2\theta_W = (m_W^2/m_Z^2)/\rho_0 = 0.776368\), so that the model
reproduces both \(m_W\) and \(m_Z\); every scale in GeV moved by
\(+8.5\times10^{-4}\) (\(m_{KK} = 121.32\) TeV, \(H = 14.296\) TeV) and
the dependent computations were rerun.

\hypertarget{the-findings-of-the-review-of-version-2-and-their-incorporation}{%
\section{The findings of the review of version 2 and their
incorporation}\label{the-findings-of-the-review-of-version-2-and-their-incorporation}}

The review agent (2026-10-18, 1858 s, 82 tool calls) read the five model
documents of version 2, compared every computed number with the outputs
(all agree except three roundings) and recomputed key results
independently (the mixing matrices, the long-baseline and reactor
probabilities, the resonances, the solar survival, the onset
temperatures, the CNB, the \(Z\) widths, the freeze-out,
\(\Delta N_{eff}\), the fast-flavour growth rates, the typed background
equations against the stored symbolic forms). It reported eight errors
(A1 to A8), minor items (B) and gaps in the self-containment (C):

\begin{itemize}
\tightlist
\item
  \textbf{A1 (the most important).} The documents compared the
  freeze-out requirement of the gauge-kinetic portal,
  \(\langle\sigma v\rangle = 2.01\times10^{-25}\) cm\(^3\)/s per
  species, with \((4\pi/M^2)\sqrt{x_f/\pi} = 2.32\times10^{-25}\)
  cm\(^3\)/s and called it 87 \% of a unitarity estimate. The operator
  annihilates only from \(^3P_0\); with the spin average the bound is
  \(\sigma \le \pi/(4p^2)\) and
  \(\langle\sigma v\rangle \le \sqrt{\pi x_f}/M^2 = 5.79\times10^{-26}\)
  cm\(^3\)/s, which the requirement exceeds 3.47 times. Incorporated:
  the computation (\texttt{um2\_analysis.py}, check B3) now computes the
  bound, its numerical thermal average, the velocity \(v_c = 0.263\)
  above which the effective cross-section exceeds it and the fraction
  (79 \%) of pairs above \(v_c\); all five documents state that the
  relic abundance of the dark matter is OPEN in version 2 (an additional
  annihilation channel or a non-thermal origin is needed), and the
  verdict on dark matter reads ``a candidate without a relic
  mechanism''.
\item
  \textbf{A2.} ``No tip operator reaches \(\Psi_B\)'' holds only for
  chiral \(P_+\) tip currents and for operators built on
  \(\bar\Psi\Psi\); corrected.
\item
  \textbf{A3.} The fast-flavour dispersion relation mixed lower- and
  upper-index shifts and did not state its metric; it is now written
  with \(\Lambda^0\), \(\Lambda^3\) and
  \(\eta^{\mu\nu} = \mathrm{diag}(1, -1)\) on \((x_4, x_3)\) in every
  document (the growth rates were confirmed).
\item
  \textbf{A4.} \(c_Z = c'_Z = 0\) was labelled DERIVED or ``required'';
  the invisible width bounds the even part and \(N_{eff}\) disfavours,
  but does not exclude, the oriented part, so it is now ASSUMED
  everywhere.
\item
  \textbf{A5.} The flavour, decoupling and fast-flavour equations wrote
  \(t\) for the time \(x_4\) and \(z\) for the axis \(x_3\); corrected
  throughout.
\item
  \textbf{A6 to A8.} Roundings: the adiabatic solar \(P_{ee}\) at 15 MeV
  is \(0.306\) (the value \(0.304\) is the numerical one); the \(\tau\)
  singlet bulk mass is \(-9.205H\); the Kohn-Sham Richardson errors are
  below \(2.2\times10^{-9}\) relative.
\item
  \textbf{B.} \(J_{max}\) computed (\(0.0335\)) against NuFIT's value
  (\(0.0333 \pm 0.0007\)) now distinguished; the conventions of the
  supernova resonances stated (two-flavour conditions); the tip matrices
  equal the mixing matrices only with \(\kappa_e = \kappa_u = 1\); the
  measured values enter as fit targets, direct inputs or comparisons;
  the single-species requirement \((17.6\ \text{TeV})^{-3}\) is a
  version-2 number; the sum rules are absolute (\(1.0\times10^{-8}\);
  relative \(7.8\times10^{-4}\) and \(6.0\times10^{-4}\) for \(W^\pm\)
  and \(Z\)); the 50-digit roots are compared with the stored masses
  (\(4.5\times10^{-11}\)); the temperature of the neutrino spectrum in
  the cosmology; the difference \(0.0044\) to \(0.0053\) from the
  precision \(N_{eff}\); the labels ESTIMATE and STANDARD; notational
  clashes (\(H\) as the Hubble rate, \(M_{Pl}\), \(\Lambda_0\),
  \(G(u)\), \(\mu_N\), the portal amplitude \(e\), \(\alpha\), \(F\) and
  \(F_h\), the field names of PROVENANCE) added to the symbol lists; the
  dark-matter mass written \(44.35\) TeV everywhere.
\item
  \textbf{C.} Equations or inputs used but not written in the same
  document (among them the relic and \(\Delta N_{eff}\) formulas, the
  Einstein-Boltzmann equations, the BAO \(\chi^2\) and data, the
  nonlinear background equations, the fermion equations with their
  boundary conditions, the decoupling equations, the CNB, Landau-Zener,
  adiabatic solar, condensate and pairing formulas, the Kohn-Sham
  equations behind the quoted equations of state, the direct-detection
  inputs) are now written out in every document that uses them.
\end{itemize}

All items were incorporated on 2026-10-18; after the change of the
analysis computation its 23 checks pass, and the total remains 205.

\end{document}
